\secciongrande{Integración del código con las etiquetas html integración: Online o en archivo externo Los navegadores aceptan la integración del código en html de diversas formas.}

\resumen{
  La integración de código JavaScript en HTML puede realizarse en línea mediante la etiqueta \texttt{<script>} o a través de un archivo externo utilizando el atributo \texttt{src}. La forma externa es la más recomendable por su modularidad, facilidad de depuración y rapidez de carga. Es crucial gestionar las rutas de los archivos (relativa, absoluta o remota) y considerar alternativas como \texttt{<noscript>} o la sintaxis \texttt{>}
}

 para xhtml.

\definicion{Integración en línea}{Inclusión directa del código JavaScript dentro de las etiquetas \texttt{<script>} del documento HTML, utilizando el atributo \texttt{type="text/javascript"}.}

\definicion{Archivo externo}{Referencia a un archivo JavaScript con extensión \texttt{.js} mediante el atributo \texttt{src} dentro de la etiqueta \texttt{<script>}.}

\subtituloA{Métodos de integración}

\lista{
  \textbf{Código en línea:} %%VERBATIM\_BLOCK\_0%% |
  \textbf{Archivo externo:} %%VERBATIM\_BLOCK\_1%%
}
